\paragraph{Proof.} For regime zero, a simple graph on \(D\) vertices has at most \(\binom D2\) adjacencies and hence at most \(2\binom D2=D(D-1)\) endpoint positions.  Each of the \(K\) nonobservational graphs has \(D+1\) vertices and therefore at most \(2\binom{D+1}{2}=D(D+1)\) endpoint positions.  Thus\[|\mathcal E(\mathbf M)|\le D(D-1)+KD(D+1)=N_{\mathrm{end}}.\]Let \(c_t\) be the number of circular endpoints after the \(t\)-th strict update.  By \(\texttt{def:monotone-closure}\), every strict update replaces at least one circle by a tail or an arrowhead, and no later update deletes that fixed mark or changes it to the incomparable fixed mark.  Consequently \(c_{t+1}\le c_t-1\).  Since \(0\le c_0\le |\mathcal E(\mathbf M)|\le N_{\mathrm{end}}\), there are at most \(N_{\mathrm{end}}\) strict endpoint updates.  The fixed scan order in \(\texttt{def:closure-runtime-certificate}\) performs a further scan after the last successful scan; that scan makes no update and certifies the least fixed point.  Hence there are at most \(N_{\mathrm{end}}+1\) exhaustive scans.

For a fixed rule \(r\), \(\texttt{def:finite-rule-encoding}\) enumerates at most \((K+1)^{a_r}\) assignments of its typed regime placeholders and at most \((D+1)^{b_r}\) assignments of its vertex placeholders.  Its finite adjacency and mark premises take constant time per assignment on the matrix encoding, while all finite-state path premises are decided by the stipulated product-graph breadth-first searches within the declared cubic envelope.  Therefore one exhaustive scan for \(r\) costs\[T_r(D,K)=O\!\left((K+1)^{a_r}(D+1)^{b_r+3}\right).\]Enumerating triples to materialize all unshielded colliders costs\[O\!\left(D^3+K(D+1)^3\right)=O\!\left((K+1)(D+1)^3\right),\]and orienting the existing indicator endpoints is of smaller order.  Summing the per-rule cost over at most \(N_{\mathrm{end}}+1\) scans gives\[T_{\mathrm{cl}}(D,K)=O\!\left((K+1)(D+1)^3+(N_{\mathrm{end}}+1)\sum_{r\in\mathsf R_{\mathrm{SI}}}T_r(D,K)\right).\]The rule set is finite and independent of \(D,K\), and each \(a_r,b_r\) is a fixed integer.  The displayed expression is therefore a polynomial in \(D,K\).  Finally, \(L_{\mathrm{in}}=D^2+K(D+1)^2\) bounds both \(D\) and \(K\) from above for \(D\ge1\), so the same expression is polynomial in the encoded input length.  This proves every asserted termination and runtime claim. \(\square\)